% =====================================================================
\section{Model 2 --- Data-Driven Model and Optimal Control}
\label{sec:model2}
% =====================================================================

Model~1 needs the physics: $m$, $J$, rotor constants. Model~2 asks a different question: \emph{if we only had
recorded flight data, could we still build a good controller?} It learns a discrete-time linear model from
data (blocks 1--2). It then designs an optimal controller on that model in two ways: LQR, which is optimal and
simple but unaware of limits (blocks 3--6), and MPC, which is optimal and respects limits but needs an optimizer
(blocks 7--9). Finally it flies the real, nonlinear drone with the result (block~10).

All matrices here are \emph{discrete-time}: $\vx_{k+1}=A\vx_k+B\vu_k$ with $\Delta t=0.04$\,s, which is the 25\,Hz
\code{control\_rate} of \code{controller.yaml}.

% ---------------------------------------------------------------------
\subsection{Block 1 --- Nonlinear flight data}
\label{sec:m2b1}
% ---------------------------------------------------------------------

\begin{eqbox}{cPlant}{Equation M2.1 \quad Sampled state and input (deviation from hover)}
\begin{equation}
  \vx_k=\begin{bmatrix}\vp(t_k)-\vp_{\mathrm{ref}}\\ \vv(t_k)\\ \Log\big(R(t_k)\big)\\ \vw(t_k)\end{bmatrix}\in\R^{12},\qquad
  \vu_k=\begin{bmatrix}T(t_k)-mg_0\\ \vtau(t_k)\end{bmatrix}\in\R^{4},\qquad t_k=k\,\Delta t . \tag{M2.1}\label{eq:m2.1}
\end{equation}
\end{eqbox}

\begin{dimtable}
\vx_k & state deviation at sample $k$ & \R^{12} & m, m/s, rad, rad/s\\
\vu_k & input deviation at sample $k$ & \R^{4} & N, N\,m\\
\vp_{\mathrm{ref}} & hover point & \R^3 & m ($[0,0,2]$)\\
\Delta t & sample time & \R & s (0.04)\\
N & number of samples in a flight & \mathbb N & -- (625 per 25\,s flight)\\
\end{dimtable}

\subsubsection*{Derivation and reasoning}

\textbf{Operating point.} Hover is an equilibrium of \eqref{eq:m1.1}--\eqref{eq:m1.4}: with $\vv=0$, $R=I$, $\vw=0$,
$T=mg_0$, $\vtau=0$, every derivative is zero. Linearizing about it gives $\delta\dot{\vx}=A_c\delta\vx+B_c\delta\vu$
with \emph{no constant term}. Measuring the state and input as deviations from hover is what lets the pure linear
form $\vx_{k+1}=A\vx_k+B\vu_k$ fit the data. If the raw thrust $T$ were used, gravity would appear as a constant
offset that this form cannot represent (correction~6 in \cref{sec:errata}).

\textbf{Attitude as a vector.} DMDc needs vectors. $\Log(R)$ turns the $3\times3$ rotation into the 3-vector rotation
angle, which near hover is the familiar (roll, pitch, yaw).

\textbf{Sampling.} The computer applies $\vu_k$ and holds it for $\Delta t$ (zero-order hold). Hence the natural model
is discrete.

\textbf{Persistent excitation.} A free quadrotor is unstable (\cref{sec:m1b3}), so data is always recorded with some
stabilising pilot in the loop, here a simple PD controller (the role ArduPilot plays). If the input were
\emph{only} feedback, $\vu_k=-K_0\vx_k$, the rows of $U$ would be linear combinations of the rows of $X$. The data could
not separate ``the drone'' from ``the pilot''. A small random excitation is therefore added on top
($\pm2$\,N thrust, $\pm0.03$\,N\,m torque, held for 0.2\,s), which makes the stacked data matrix full rank.

\begin{toybox}[the recorded data set]
Five flights of 25\,s at 25\,Hz, each starting from a random point near $[0,0,2]$\,m: $5\times625$ samples of
$\vx_k\in\R^{12}$ and $\vu_k\in\R^4$, with small logging noise. Flights 1--4 are used for identification and
flight~5 is kept for validation. The largest roll/pitch in the data is $6.5^\circ$, well inside the
near-linear range of \cref{sec:m1b2}.
\end{toybox}

\begin{figure}[h]
  \centering
  \includegraphics[width=0.92\linewidth]{step1_flight_data.png}
  \caption{Toy model, Model~2 step~1: example time histories and the 3D tracks of the five recorded flights.}
\end{figure}

\begin{projbox}
In the project, this block is a \code{ros2 bag record} of \code{/ap/v1/pose/filtered},
\code{/ap/v1/twist/filtered} and the commands we send, recorded from ArduPilot SITL in Gazebo or from the real
drone. Nothing about mass or inertia needs to be known.
\end{projbox}

\begin{codebox}
\intoy \code{model2\_dmdc\_lqr\_mpc/step1\_flight\_data.m}, \code{lib/state\_to\_dev.m}, \code{lib/baseline\_pilot.m}.
\end{codebox}

% ---------------------------------------------------------------------
\subsection{Block 2 --- DMDc: Dynamic Mode Decomposition with control}
\label{sec:m2b2}
% ---------------------------------------------------------------------

\begin{eqbox}{cOuter}{Equations M2.2 -- M2.4 \quad Snapshot matrices and least-squares model}
\begin{align}
  X=\big[\vx_1\ \cdots\ \vx_{N-1}\big],\quad X^+&=\big[\vx_2\ \cdots\ \vx_{N}\big],\quad U=\big[\vu_1\ \cdots\ \vu_{N-1}\big] \tag{M2.2}\label{eq:m2.2}\\
  \big[\,A\ \ B\,\big]&=X^+\begin{bmatrix}X\\ U\end{bmatrix}^{\dagger} \tag{M2.3}\label{eq:m2.3}\\
  \vx_{k+1}&\approx A\,\vx_k+B\,\vu_k \tag{M2.4}\label{eq:m2.4}
\end{align}
\end{eqbox}

\begin{dimtable}
X,\ X^+ & state snapshots and the same shifted by one sample & \R^{12\times(N-1)} & mixed\\
U & input snapshots & \R^{4\times(N-1)} & N, N\,m\\
\Omega=[X;U] & stacked regressor matrix & \R^{16\times(N-1)} & mixed\\
\Omega^\dagger & pseudo-inverse of $\Omega$ & \R^{(N-1)\times16} & mixed\\
A & learned discrete state matrix & \R^{12\times12} & --\\
B & learned discrete input matrix & \R^{12\times4} & mixed\\
\end{dimtable}

\subsubsection*{Derivation}

We want the pair $G=[A\ B]$ for which $\vx_{k+1}=A\vx_k+B\vu_k=G\,[\vx_k;\vu_k]$ holds as well as possible for
\emph{every} recorded sample. Written for all samples at once, this is $X^+\approx G\,\Omega$. The least-squares choice
minimises the sum of squared one-step prediction errors:
\[
  \min_G\ \big\|X^+-G\,\Omega\big\|_F^2 .
\]
Setting the gradient to zero, $-2\,(X^+-G\Omega)\Omega^\top=0$, gives the \emph{normal equations}
$G\,\Omega\Omega^\top=X^+\Omega^\top$. If $\Omega$ has full row rank (16, guaranteed by the excitation), the
$16\times16$ matrix $\Omega\Omega^\top$ is invertible and
\[
  G=X^+\Omega^\top\big(\Omega\Omega^\top\big)^{-1}=X^+\Omega^\dagger .
\]
In practice $\Omega^\dagger$ is computed from the SVD $\Omega=U_\Omega\Sigma V_\Omega^\top$ as
$V_\Omega\Sigma^{-1}U_\Omega^\top$, and tiny singular values can be truncated to suppress noise
(Proctor, Brunton \& Kutz, 2016). Splitting the columns of $G$ gives $A$ (first 12) and $B$ (last 4).

\textbf{Relation to Model~1.} If the data came from the linearized physics, the exact answer would be the
zero-order-hold discretization of \eqref{eq:AB}:
$A=e^{A_c\Delta t}$, $B=\int_0^{\Delta t}e^{A_cs}ds\,B_c$. DMDc recovers these numbers from data alone.

\begin{toybox}[a scalar system by hand]
True system $x_{k+1}=0.9x_k+0.5u_k$. Three recorded steps:
$X=[1,\ 1.4,\ 0.76]$, $X^+=[1.4,\ 0.76,\ 0.934]$, $U=[1,\ -1,\ 0.5]$.
\[
\Omega\Omega^\top=\begin{bmatrix}3.5376&-0.02\\-0.02&2.25\end{bmatrix},\qquad
X^+\Omega^\top=\begin{bmatrix}3.17384&1.107\end{bmatrix},
\]
\[
[a\ \ b]=\frac{1}{7.9592}\begin{bmatrix}3.17384&1.107\end{bmatrix}\begin{bmatrix}2.25&0.02\\0.02&3.5376\end{bmatrix}
=\begin{bmatrix}0.9000&0.5000\end{bmatrix}.
\]
The true model is recovered exactly, because the data is noise-free.

\textbf{The drone} (2500 snapshots, rank$[X;U]=16$):
\begin{center}\small
\begin{tabular}{@{}llrr@{}}
\toprule
entry & physical meaning & DMDc & physics (ZOH)\\
\midrule
$A_{1,4}$ & $\Delta t$ (position from velocity) & 0.0399 & 0.0400\\
$A_{4,8}$ & $\approx g_0\Delta t$ ($v_x$ from pitch) & 0.3909 & 0.3924\\
$A_{5,7}$ & $\approx-g_0\Delta t$ ($v_y$ from roll) & $-0.3862$ & $-0.3924$\\
$B_{6,1}$ & $\Delta t/m$ ($v_z$ from thrust) & 0.0268 & 0.0267\\
$B_{10,2}$ & $\Delta t/J_x$ (roll rate from roll torque) & 2.0007 & 2.0000\\
$B_{12,4}$ & $\Delta t/J_z$ (yaw rate from yaw torque) & 1.0096 & 1.0000\\
\bottomrule
\end{tabular}
\end{center}
Overall $\|A-A_{\text{phys}}\|/\|A_{\text{phys}}\|=0.95\%$ and $\|B-B_{\text{phys}}\|/\|B_{\text{phys}}\|=0.81\%$. On the unseen
flight, the 1-second-ahead position prediction error is 1.21\,cm, compared with 1.14\,cm for the exact physics
model. The learned mass is $\Delta t/B_{6,1}=1.49$\,kg and the learned roll inertia is $\Delta t/B_{10,2}=0.0200$\,kg\,m$^2$.
\end{toybox}

\begin{figure}[h]
  \centering
  \includegraphics[width=0.9\linewidth]{step2_dmdc_validation.png}
  \caption{Toy model, Model~2 step~2. Left: eigenvalues of the learned and physical $A$, all at or near $1$
  (integrators). Right: one-step predictions on a flight that was \emph{not} used for learning.}
\end{figure}

\begin{projbox}
DMDc is ``curve fitting for dynamics''. It looks at thousands of moments of the form ``the drone was here, we
pushed like this, and 40\,ms later it was there'', and finds the single linear rule that explains them best. For our
drone this means we can re-identify the model whenever something changes (a heavier camera, a new battery, a
different frame) simply by flying for two minutes.
\end{projbox}

\begin{codebox}
\intoy \code{step2\_dmdc.m}, \code{lib/dmdc.m}, \code{lib/c2d\_zoh.m}.
\end{codebox}

% ---------------------------------------------------------------------
\subsection{Block 3 --- Bellman equation (optimal control)}
\label{sec:m2b3}
% ---------------------------------------------------------------------

\begin{eqbox}{cGuid}{Equations M2.5 -- M2.7 \quad Cost, Bellman equation, quadratic value function}
\begin{align}
  J&=\sum_{k=0}^{\infty}\big(\vx_k^\top Q\,\vx_k+\vu_k^\top R\,\vu_k\big) \tag{M2.5}\label{eq:m2.5}\\
  V(\vx_k)&=\min_{\vu_k}\Big[\vx_k^\top Q\vx_k+\vu_k^\top R\vu_k+V(\vx_{k+1})\Big],\qquad \vx_{k+1}=A\vx_k+B\vu_k \tag{M2.6}\label{eq:m2.6}\\
  V(\vx)&=\vx^\top P\,\vx \tag{M2.7}\label{eq:m2.7}
\end{align}
\end{eqbox}

\begin{dimtable}
J & total cost of a flight & \R & -- (dimensionless)\\
Q & state weight, $Q=Q^\top\succeq0$ & \R^{12\times12} & 1/(unit of $x_i$)$^2$\\
R & input weight, $R=R^\top\succ0$ & \R^{4\times4} & 1/N$^2$, 1/(N\,m)$^2$\\
V(\vx) & value function: lowest possible cost from state $\vx$ & \R & --\\
P & value matrix, $P=P^\top\succeq0$ & \R^{12\times12} & as $Q$\\
\end{dimtable}

\subsubsection*{Derivation}

\textbf{(M2.5) The cost.} $\vx^\top Q\vx$ is a weighted sum of squared errors and $\vu^\top R\vu$ a weighted sum of
squared effort. Minimising $J$ trades accuracy against effort. We choose the weights by \emph{Bryson's rule},
$Q_{ii}=1/x_{i,\max}^2$ and $R_{jj}=1/u_{j,\max}^2$, where $x_{i,\max}$ is the largest acceptable deviation. A 50\,cm
position error then costs as much as a 1\,m/s speed error, a $0.2$\,rad tilt, or 5\,N of extra thrust.

\textbf{(M2.6) Principle of optimality} (Bellman). Split the infinite sum into the first step and the rest:
\[
  V(\vx_k)=\min_{\vu_k,\vu_{k+1},\dots}\Big[\ell(\vx_k,\vu_k)+\sum_{j>k}\ell(\vx_j,\vu_j)\Big]
  =\min_{\vu_k}\Big[\ell(\vx_k,\vu_k)+\underbrace{\min_{\vu_{k+1},\dots}\sum_{j>k}\ell(\vx_j,\vu_j)}_{V(\vx_{k+1})}\Big].
\]
Whatever the first move is, the remaining moves must be optimal for the state it leads to. Otherwise we could
improve them and lower the total. The infinite-dimensional problem collapses to a one-step problem.

\textbf{(M2.7) Quadratic value function.} Suppose $V(\vx_{k+1})=\vx_{k+1}^\top P\vx_{k+1}$. The bracket in \eqref{eq:m2.6}
is then a quadratic function of $\vu_k$; its minimiser is linear in $\vx_k$ (block~4), and the minimum value is again a
quadratic form in $\vx_k$ (block~5). The quadratic shape is preserved by the Bellman equation, so the guess
$V=\vx^\top P\vx$ is self-consistent, and all that remains is to find $P$.

\begin{toybox}[one Bellman step by hand]
$x_{k+1}=1.1x_k+0.5u_k$ (unstable: grows 10\% per step), $q=r=1$, next-step value $V(x)=1\cdot x^2$, current $x=1$:
\[
  \text{RHS}(u)=1+u^2+(1.1+0.5u)^2,\qquad \frac{d}{du}=2u+(1.1+0.5u)=0\ \Rightarrow\ u^*=-\frac{1.1}{2.5}=-0.44,
\]
\[
  V(1)=1+0.1936+(1.1-0.22)^2=1.968 .
\]
\code{fminsearch} on the same function returns $u^*=-0.4400$ and $V=1.9680$.

\textbf{Drone} (DMDc model, $V_{\text{next}}=\vx^\top Q\vx$, drone 50/30/20\,cm off target): brute-force minimisation over
$\vu\in\R^4$ and the closed form both give $\vu^*=[-0.0104,-0.0001,0.0001,0.0000]$ and the same value 3.0393.
Doing nothing ($\vu=0$) costs 14\,593 after 10\,s and keeps growing.
\end{toybox}

\begin{projbox}
$Q$ and $R$ are where \emph{we} write down what ``good flying'' means for our mission. For mapping with a camera we
care about smooth, gentle motion (blur-free images), so tilt and rate errors are weighted heavily and speed is
limited to the \code{max\_velocity} of 1\,m/s. The Bellman equation turns that wish list into a mathematical problem with a
unique best answer.
\end{projbox}

\begin{codebox}
\intoy \code{step3\_bellman\_cost.m}, \code{lib/bellman\_backup.m}.
\end{codebox}

% ---------------------------------------------------------------------
\subsection{Block 4 --- Linear Quadratic Regulator (LQR)}
\label{sec:m2b4}
% ---------------------------------------------------------------------

\begin{eqbox}{cGuid}{Equations M2.8 -- M2.9 \quad Optimal gain and control law}
\begin{align}
  K&=\big(R+B^\top PB\big)^{-1}B^\top PA \tag{M2.8}\label{eq:m2.8}\\
  \vu_k&=-K\,\vx_k \tag{M2.9}\label{eq:m2.9}
\end{align}
\end{eqbox}

\begin{dimtable}
K & state-feedback gain & \R^{4\times12} & N/m, N\,s/m, N\,m/rad, \dots\\
R+B^\top PB & curvature of the one-step cost in $\vu$ & \R^{4\times4} & --\\
\end{dimtable}

\subsubsection*{Derivation}

Insert $V(\vx_{k+1})=(A\vx+B\vu)^\top P(A\vx+B\vu)$ into the bracket of \eqref{eq:m2.6} and differentiate with respect to
$\vu$ (using $\partial(\vu^\top M\vu)/\partial\vu=2M\vu$ for symmetric $M$):
\[
  \frac{\partial}{\partial\vu}\Big[\vx^\top Q\vx+\vu^\top R\vu+(A\vx+B\vu)^\top P(A\vx+B\vu)\Big]
  =2R\vu+2B^\top P(A\vx+B\vu)=\bm0 .
\]
Collecting the $\vu$ terms: $(R+B^\top PB)\,\vu=-B^\top PA\,\vx$. Because $R\succ0$ and $P\succeq0$, the matrix
$R+B^\top PB$ is positive definite. It is therefore invertible, and its positive curvature (the Hessian $2(R+B^\top PB)$)
guarantees a minimum. Hence $\vu^*=-(R+B^\top PB)^{-1}B^\top PA\,\vx=-K\vx$.

\begin{toybox}[the scalar system]
With the converged $P=3.1215$ (block~6): $K=\dfrac{ab\,P}{r+b^2P}=\dfrac{1.1\cdot0.5\cdot3.1215}{1+0.25\cdot3.1215}=0.9643$.
Closed loop: $x_{k+1}=(a-bK)x_k=0.6179\,x_k$. Without control the state grows 10\% per step; with LQR it shrinks
38\% per step.

\textbf{Drone gain} (from block~6), two of its 48 numbers:
\[
\delta T=-9.04\,\delta z-6.92\,v_z,\qquad
\tau_x=0.273\,\delta y+0.321\,v_y-1.444\,\theta_x-0.276\,\omega_x .
\]
If the drone is too high, thrust drops. If it is to the left ($\delta y>0$), it rolls right ($+\tau_x$, which makes $\dot v_y<0$),
but less so if it is already rolled or rolling.

With only a one-step look-ahead ($P=$ one Bellman step from $Q$), the gain gives a spectral radius
$\rho(A-BK)=1.0024>1$: \emph{unstable}. The full infinite-horizon $P$ is essential.
\end{toybox}

\begin{projbox}
The LQR controller is just a table: a $4\times12$ matrix $K$. Every 40\,ms it multiplies the 12 measured errors by $K$
and obtains thrust and three torques. It is extremely cheap to run on the drone and, for the cost we chose,
provably the best possible linear controller. Its weakness is that it knows nothing about limits.
\end{projbox}

\begin{codebox}
\intoy \code{step4\_lqr\_gain.m}, \code{lib/lqr\_gain.m}.
\end{codebox}

% ---------------------------------------------------------------------
\subsection{Block 5 --- Riccati recursion}
\label{sec:m2b5}
% ---------------------------------------------------------------------

\begin{eqbox}{cGuid}{Equation M2.10 \quad Discrete Riccati difference equation}
\begin{equation}
  P_{k+1}=A^\top P_kA-A^\top P_kB\big(R+B^\top P_kB\big)^{-1}B^\top P_kA+Q,\qquad P_0=Q\ \text{(or any terminal weight)} \tag{M2.10}\label{eq:m2.10}
\end{equation}
\end{eqbox}

\begin{dimtable}
P_k & value matrix of a problem with $k$ steps to go & \R^{12\times12} & as $Q$\\
k & iteration index (steps-to-go, \emph{not} time) & \mathbb N & --\\
\end{dimtable}

\subsubsection*{Derivation}

Substitute the optimal input $\vu^*=-K\vx$ back into the bracket of \eqref{eq:m2.6}:
\[
  V_{\text{new}}(\vx)=\vx^\top\Big[Q+K^\top RK+(A-BK)^\top P(A-BK)\Big]\vx .
\]
Expand $(A-BK)^\top P(A-BK)=A^\top PA-A^\top PBK-K^\top B^\top PA+K^\top B^\top PBK$ and add $K^\top RK$:
\[
\begin{aligned}
  &K^\top(R+B^\top PB)K=K^\top B^\top PA\quad\text{(by \eqref{eq:m2.8})}\\
  &\Longrightarrow\ V_{\text{new}}=\vx^\top\big[Q+A^\top PA-A^\top PBK\big]\vx
   =\vx^\top\big[Q+A^\top PA-A^\top PB(R+B^\top PB)^{-1}B^\top PA\big]\vx .
\end{aligned}
\]
This is \eqref{eq:m2.10}. Starting from the terminal cost $P_0$, each application gives the value matrix of a
problem that is one step longer. If $(A,B)$ is stabilisable and $Q\succ0$, the sequence converges to the
infinite-horizon matrix $P_\infty$.

\begin{toybox}[scalar recursion, $a=1.1$, $b=0.5$, $q=r=1$, $P_0=1$]
\[
P_{k+1}=1+1.21P_k-\frac{(0.55P_k)^2}{1+0.25P_k}
\]
\begin{center}\small
\begin{tabular}{l*{10}{c}}
\toprule
$k$ & 0 & 1 & 2 & 3 & 4 & 5 & 6 & 7 & 8 & $\infty$\\
$P_k$ & 1 & 1.968 & 2.596 & 2.905 & 3.036 & 3.089 & 3.109 & 3.117 & 3.120 & 3.1215\\
\bottomrule
\end{tabular}
\end{center}
(35 iterations to $10^{-14}$). For the drone, starting from $P_0=Q$, the relative change falls below $10^{-12}$
after 236 iterations, and $P_\infty$ is positive definite ($\lambda_{\min}=3.27$).
\end{toybox}

\begin{figure}[h]
  \centering
  \includegraphics[width=0.85\linewidth]{step5_riccati_convergence.png}
  \caption{Toy model, Model~2 step~5: convergence of the Riccati recursion for the learned drone model.}
\end{figure}

\begin{projbox}
Each iteration lets the controller ``look one step further into the future''. After a few hundred steps, which is 10\,s
of flight at 25\,Hz, looking further no longer changes anything. The drone then plans as if the flight never ends,
which is what an infinite-horizon LQR is.
\end{projbox}

\begin{codebox}
\intoy \code{step5\_riccati\_recursion.m}, \code{lib/riccati\_recursion.m}.
\end{codebox}

% ---------------------------------------------------------------------
\subsection{Block 6 --- DARE: Discrete Algebraic Riccati Equation}
\label{sec:m2b6}
% ---------------------------------------------------------------------

\begin{eqbox}{cGuid}{Equation M2.11 \quad Steady-state Riccati equation}
\begin{equation}
  P=A^\top PA-A^\top PB\big(R+B^\top PB\big)^{-1}B^\top PA+Q \tag{M2.11}\label{eq:m2.11}
\end{equation}
\end{eqbox}

\begin{dimtable}
P & stabilising solution, $P=P^\top\succ0$ & \R^{12\times12} & as $Q$\\
\end{dimtable}

\subsubsection*{Derivation}

If $P_k\to P$, both sides of \eqref{eq:m2.10} tend to the same limit, which gives \eqref{eq:m2.11}. Two further facts make
$P$ useful.

\textbf{Optimal cost.} Using the rearranged form of \eqref{eq:m2.11},
$P=Q+K^\top RK+(A-BK)^\top P(A-BK)$ (from the derivation of block~5), along the closed loop $\vx_{k+1}=(A-BK)\vx_k$:
\[
  \vx_k^\top P\vx_k-\vx_{k+1}^\top P\vx_{k+1}=\vx_k^\top Q\vx_k+\vu_k^\top R\vu_k .
\]
Summing from $k=0$ to $\infty$ (a telescoping sum, and $\vx_k\to0$) gives $J^*=\vx_0^\top P\vx_0$.

\textbf{Direct solution.} Instead of iterating, the optimality conditions $\vu_k=-R^{-1}B^\top\bm\lambda_{k+1}$,
$\bm\lambda_k=Q\vx_k+A^\top\bm\lambda_{k+1}$ (co-state $\bm\lambda_k=P\vx_k$) give the linear pencil
\[
  \begin{bmatrix}I&BR^{-1}B^\top\\0&A^\top\end{bmatrix}\begin{bmatrix}\vx_{k+1}\\ \bm\lambda_{k+1}\end{bmatrix}
  =\begin{bmatrix}A&0\\-Q&I\end{bmatrix}\begin{bmatrix}\vx_k\\ \bm\lambda_k\end{bmatrix}.
\]
Its 12 stable generalised eigenvectors $[X_1;X_2]$ ($|\lambda|<1$) span $\{\bm\lambda=P\vx\}$, so $P=X_2X_1^{-1}$. This is
what \code{lib/dare\_solve.m} implements, so no toolbox is needed.

\begin{toybox}[scalar DARE is a quadratic equation]
$P=1+1.21P-\dfrac{0.3025P^2}{1+0.25P}$. Multiplying by $(1+0.25P)$ and simplifying:
\[
  0.25P^2-0.46P-1=0\ \Longrightarrow\ P=\frac{0.46+\sqrt{0.46^2+1}}{0.5}=3.1215
\]
(the positive root), the limit of the recursion above. The eigenvector solver returns the same value.

\textbf{Drone:} DARE residual $4.2\times10^{-14}$ (relative); agreement with the 236-step recursion $1.1\times10^{-11}$;
closed-loop spectral radius $\rho(A-BK)=0.950<1$, so it is stable. From $\vx_0=[0.5,-0.3,0.2,0,\dots]$: the simulated cost
$\sum(\vx^\top Q\vx+\vu^\top R\vu)=41.4373$ and the predicted $\vx_0^\top P\vx_0=41.4373$ agree.
\end{toybox}

\begin{projbox}
The DARE gives the final $P$ in one shot, and from it the gain $K$ that we would upload to the drone. $P$ is also
reused: it becomes the terminal cost of the MPC (dashed arrow in \cref{fig:flow2}). That is the bridge between the two
branches.
\end{projbox}

\begin{codebox}
\intoy \code{step6\_dare.m}, \code{lib/dare\_solve.m}.
\end{codebox}

% ---------------------------------------------------------------------
\subsection{Block 7 --- Model Predictive Control (MPC)}
\label{sec:m2b7}
% ---------------------------------------------------------------------

\begin{eqbox}{cInner}{Equation M2.12 \quad Finite-horizon constrained optimal control problem}
\begin{equation}
\begin{aligned}
  \min_{\vu_0,\dots,\vu_{N-1}}\ \ &\sum_{k=0}^{N-1}\big(\vx_k^\top Q\vx_k+\vu_k^\top R\vu_k\big)+\vx_N^\top P_f\,\vx_N\\
  \text{subject to}\ \ &\vx_{k+1}=A\vx_k+B\vu_k,\quad \vx_0=\vx(t)\ \text{(measured)},\\
  &\vu_{\min}\le\vu_k\le\vu_{\max},\qquad |C\vx_k|\le\vx_{\max}\quad(k=1..N)
\end{aligned}\tag{M2.12}\label{eq:m2.12}
\end{equation}
\end{eqbox}

\begin{dimtable}
N & prediction horizon & \mathbb N & steps (25 = 1\,s)\\
P_f & terminal weight (here $P_f=P$ from the DARE) & \R^{12\times12} & as $Q$\\
\vu_{\min},\vu_{\max} & input limits (deviation) & \R^4 & N, N\,m\\
C & selects the limited states ($v_x,v_y,v_z$, roll, pitch) & \R^{5\times12} & --\\
\vx_{\max} & state limits & \R^5 & m/s, rad\\
\end{dimtable}

\subsubsection*{Derivation and reasoning}

LQR solves \eqref{eq:m2.5} without limits. Real drones have limits: thrust between roughly 5 and 25\,N, torque that the
rotors can actually produce, the \code{max\_velocity} of 1\,m/s, and a maximum tilt for a camera mount. Adding
constraints destroys the closed-form solution, so MPC solves a \emph{finite} version numerically, from the current
measured state:
\begin{itemize}
  \item the sum stops after $N$ steps, which keeps the problem finite;
  \item the infinite tail $\sum_{k\ge N}$ is replaced by the terminal cost $\vx_N^\top P_f\vx_N$.
\end{itemize}

\textbf{Why $P_f=P$ from the DARE.} By \cref{sec:m2b6}, $\vx_N^\top P\vx_N$ is \emph{exactly} the optimal cost of
the rest of the flight if no constraint is active after step $N$. With this choice, whenever the constraints are not
active, the MPC solution coincides with the LQR solution. Block~8 verifies this numerically.

\begin{toybox}[a two-step problem]
$a=1.1$, $b=0.5$, $q=r=1$, $N=2$, $P_f=3.1215$, $x_0=1$, $|u_k|\le0.5$:
\[
  \min_{u_0,u_1}\ x_0^2+u_0^2+x_1^2+u_1^2+3.1215\,x_2^2\quad\text{s.t.}\ x_{k+1}=1.1x_k+0.5u_k .
\]
\textbf{Drone:} $N=25$ (1\,s look-ahead); $5\le T\le25$\,N, i.e.\ $\delta T\in[-9.715,\,10.285]$\,N; $|\tau_{x,y}|\le0.3$,
$|\tau_z|\le0.05$\,N\,m; $|v_i|\le1$\,m/s; $|\theta_{x,y}|\le0.35$\,rad (20$^\circ$).
\end{toybox}

\begin{projbox}
MPC is a chess player. Every 40\,ms it imagines the next second of flight, finds the best sequence of moves that never
breaks a rule (``never faster than 1\,m/s'', ``never ask a motor for more than it can give''), plays only the first move,
and thinks again. For camera-based mapping, the speed and tilt limits protect image quality.
\end{projbox}

\begin{codebox}
\intoy \code{step7\_mpc\_problem.m}, \code{lib/mpc\_build.m}.
\end{codebox}

% ---------------------------------------------------------------------
\subsection{Block 8 --- Optimization formulation: Quadratic Program}
\label{sec:m2b8}
% ---------------------------------------------------------------------

\begin{eqbox}{cInner}{Equations M2.13 -- M2.15 \quad Stacked prediction, cost and constraints}
\begin{align}
  \mathbf X&=S_x\vx_0+S_u\mathbf U,\qquad
  \mathbf X=\begin{bmatrix}\vx_1\\ \vdots\\ \vx_N\end{bmatrix},\\
  \mathbf U=\begin{bmatrix}\vu_0\\ \vdots\\ \vu_{N-1}\end{bmatrix},\\
  S_x=\begin{bmatrix}A\\ A^2\\ \vdots\\ A^N\end{bmatrix},\\
  S_u=\begin{bmatrix}B&0&\cdots&0\\ AB&B&\cdots&0\\ \vdots&&\ddots&\vdots\\ A^{N-1}B&\cdots&AB&B\end{bmatrix} \tag{M2.13}\label{eq:m2.13}\\[4pt]
  &\min_{\mathbf U}\ \tfrac12\mathbf U^\top H\mathbf U+\bm f^\top\mathbf U,\quad
  H=2\big(S_u^\top\bar QS_u+\bar R\big),\ \ \bm f=2S_u^\top\bar QS_x\vx_0 \tag{M2.14}\label{eq:m2.14}\\[2pt]
  &\text{s.t.}\quad G\mathbf U\le\bm h,\qquad A_{eq}\mathbf U=\bm b_{eq} \tag{M2.15}\label{eq:m2.15}
\end{align}
with $\bar Q=\diag(Q,\dots,Q,P_f)$ and $\bar R=\diag(R,\dots,R)$.
\end{eqbox}

\begin{dimtable}
\mathbf X,\ \mathbf U & stacked predicted states / inputs & \R^{12N},\ \R^{4N} & mixed\\
S_x,\ S_u & prediction matrices & \R^{12N\times12},\ \R^{12N\times4N} & --\\
\bar Q,\ \bar R & block-diagonal weights & \R^{12N\times12N},\ \R^{4N\times4N} & --\\
H,\ \bm f & QP Hessian and linear term & \R^{4N\times4N},\ \R^{4N} & --\\
G,\ \bm h & inequality constraints & \R^{n_c\times4N},\ \R^{n_c} & mixed\\
A_{eq},\ \bm b_{eq} & equality constraints (optional) & \R^{n_e\times4N},\ \R^{n_e} & mixed\\
\end{dimtable}

\subsubsection*{Derivation}

\textbf{(M2.13) Prediction.} Apply $\vx_{k+1}=A\vx_k+B\vu_k$ repeatedly:
$\vx_1=A\vx_0+B\vu_0$, $\vx_2=A^2\vx_0+AB\vu_0+B\vu_1$, and in general
\[
  \vx_k=A^k\vx_0+\sum_{j=0}^{k-1}A^{k-1-j}B\vu_j .
\]
Stacking $k=1..N$ gives \eqref{eq:m2.13}. The future states are \emph{linear} in the unknown inputs.

\textbf{(M2.14) Cost.} The cost of \eqref{eq:m2.12} is $\vx_0^\top Q\vx_0+\mathbf X^\top\bar Q\mathbf X+\mathbf U^\top\bar R\mathbf U$.
Substitute $\mathbf X$:
\[
  \mathbf U^\top\big(S_u^\top\bar QS_u+\bar R\big)\mathbf U+2\,\vx_0^\top S_x^\top\bar QS_u\,\mathbf U+\underbrace{\vx_0^\top(Q+S_x^\top\bar QS_x)\vx_0}_{\text{does not depend on }\mathbf U}.
\]
Dropping the constant and writing it as $\tfrac12\mathbf U^\top H\mathbf U+\bm f^\top\mathbf U$ gives $H$ and $\bm f$ in
\eqref{eq:m2.14}. $H\succ0$ because $\bar R\succ0$, so the problem is \emph{convex}, with a single global minimum
and no local traps.

\textbf{(M2.15) Constraints.} Input limits: $[I;-I]\,\mathbf U\le[\mathbf U_{\max};-\mathbf U_{\min}]$. State limits
$|\bar C\mathbf X|\le\bar{\vx}_{\max}$ become
$\bar CS_u\mathbf U\le\bar{\vx}_{\max}-\bar CS_x\vx_0$ and $-\bar CS_u\mathbf U\le\bar{\vx}_{\max}+\bar CS_x\vx_0$.
Stacking everything gives $G\mathbf U\le\bm h(\vx_0)$. Equality constraints appear, for example, as a terminal
condition $\vx_N=0$, i.e.\ (last block row of $S_u$)$\,\mathbf U=-A^N\vx_0$, or in the ``sparse'' formulation, where both
$\mathbf X$ and $\mathbf U$ are unknowns and the dynamics themselves are the equalities.

\textbf{Solver.} The toy model solves the QP with a primal--dual interior-point method (\code{lib/qp\_solve.m}).
It applies Newton steps on the KKT conditions $H\mathbf U+\bm f+G^\top\bm z+A_{eq}^\top\bm y=0$, $G\mathbf U+\bm s=\bm h$,
$s_iz_i=0$, keeping $\bm s,\bm z>0$. It replaces \code{quadprog}, so no Optimization Toolbox is needed.

\begin{toybox}[the two-step problem of block 7, solved by hand]
\[
S_x=\begin{bmatrix}1.1\\1.21\end{bmatrix},\quad S_u=\begin{bmatrix}0.5&0\\0.55&0.5\end{bmatrix},\quad
\bar Q=\diag(1,\,3.1215),
\]
\[
H=2\left(\begin{bmatrix}1.1942&0.8584\\0.8584&0.7804\end{bmatrix}+I\right)=\begin{bmatrix}4.3885&1.7168\\1.7168&3.5607\end{bmatrix},\qquad
\bm f=\begin{bmatrix}5.2547\\3.7770\end{bmatrix}.
\]
\begin{itemize}
\item \textbf{No limits:} $\mathbf U^*=-H^{-1}\bm f=[-0.9643,\ -0.5958]$. The first move equals the LQR move
      $-Kx_0=-0.9643$ \emph{exactly}, as promised by $P_f=P$.
\item \textbf{With $|u|\le0.5$:} $\mathbf U^*=[-0.5,\ -0.5]$. Both limits are active and the KKT multipliers are positive.
\item \textbf{With $A_{eq}\mathbf U=b_{eq}$ forcing $x_2=0$} ($[0.55\ \ 0.5]\,\mathbf U=-1.21$): $\mathbf U^*=[-1.3057,\ -0.9837]$.
      Check: $0.55(-1.3057)+0.5(-0.9837)=-1.21$.
\end{itemize}
\textbf{Drone ($N=25$):} 100 decision variables and 450 inequality constraints. For a 5\,cm offset (no limit active) the
MPC input equals $-K\vx_0$ to all printed digits. For a 2\,m offset the unconstrained LQR would demand 1.39 times the
allowed input; the QP saturates the thrust at $+10.285$\,N and plans a velocity that touches but never exceeds
1\,m/s. Solved in 10 interior-point iterations (about 8\,ms).
\end{toybox}

\begin{projbox}
Here the physics disappears into a few matrices. To the solver, ``fly the drone well without breaking the
rules'' is just ``find the lowest point of a bowl ($H$, $\bm f$) inside a fence ($G$, $\bm h$)''. Because the bowl has
one bottom, the answer is unique and can be found reliably in milliseconds. That is what makes MPC feasible at 25\,Hz.
\end{projbox}

\begin{codebox}
\intoy \code{step8\_qp\_formulation.m}, \code{lib/mpc\_build.m}, \code{lib/qp\_solve.m}.
\end{codebox}

% ---------------------------------------------------------------------
\subsection{Block 9 --- MPC control input (receding horizon)}
\label{sec:m2b9}
% ---------------------------------------------------------------------

\begin{eqbox}{cInner}{Equation M2.16 \quad Receding-horizon control law}
\begin{equation}
  \mathbf U^*(\vx_k)=\big[\vu_0^*,\ \vu_1^*,\ \dots,\ \vu_{N-1}^*\big]=\arg\min\text{QP}(\vx_k)
  \quad\Longrightarrow\quad \vu_k=\vu_0^*(\vx_k),\quad k\leftarrow k+1 \tag{M2.16}\label{eq:m2.16}
\end{equation}
\end{eqbox}

\begin{dimtable}
\mathbf U^*(\vx_k) & optimal plan computed at time $k$ & \R^{4N} & N, N\,m\\
\vu_0^* & first element of the plan, the only one applied & \R^{4} & N, N\,m\\
\end{dimtable}

\subsubsection*{Derivation and reasoning}

The plan $\mathbf U^*$ is optimal \emph{for the model}. The model is never perfect (it was learned, and wind is not in
it). Re-solving at every sample from the newly measured state $\vx_{k+1}$ turns the open-loop plan into a feedback law
$\vu_k=\kappa(\vx_k)$: errors are corrected 25 times per second instead of accumulating. Applying $\vu_0^*$ and shifting the
horizon forward by one step is the \emph{receding horizon} principle.

\begin{toybox}[8 s on the learned model]
Start 1.5\,m, 1.0\,m and 1.2\,m away from the setpoint. At $k=0$ the plan begins with
$\vu_0^*=[10.285,\,-0.279,\,-0.300,\,0.002]$: full upward thrust and the maximum allowed torque.
Over the run, MPC respects $|v_i|\le1$\,m/s (maximum 1.000), while LQR clipped to the same input limits reaches 1.172\,m/s.
MPC converges to $\|\vx\|=1.3\times10^{-5}$. The state limits never had to be relaxed. The mean solve time is 2.2\,ms
against a 40\,ms budget.
\end{toybox}

\begin{figure}[h]
  \centering
  \includegraphics[width=0.9\linewidth]{step9_receding_horizon.png}
  \caption{Toy model, Model~2 step~9: receding-horizon MPC versus input-clipped LQR on the learned model.}
\end{figure}

\begin{projbox}
``Plan one second ahead, act for 40 milliseconds, look again.'' Because the plan is recomputed from fresh EKF
estimates, the MPC inherits the robustness of feedback while still respecting every limit.
\end{projbox}

\begin{codebox}
\intoy \code{step9\_mpc\_control\_input.m}, \code{lib/mpc\_solve.m}.
\end{codebox}

% ---------------------------------------------------------------------
\subsection{Block 10 --- Apply to the drone}
\label{sec:m2b10}
% ---------------------------------------------------------------------

\begin{eqbox}{cPlant}{Equations M2.17 -- M2.18 \quad From optimal input to motor speeds}
\begin{align}
  \vu_k^{\text{abs}}&=\vu_{\text{trim}}+\begin{cases}\vu_0^*(\vx_k)&\text{(MPC)}\\ -K\vx_k&\text{(LQR)}\end{cases}
   =\begin{bmatrix}T_k\\ \vtau_k\end{bmatrix},\qquad \vu_{\text{trim}}=\begin{bmatrix}mg_0\\ \bm0\end{bmatrix} \tag{M2.17}\label{eq:m2.17}\\
  \vf&=\Gamma^{-1}\vu_k^{\text{abs}},\qquad \Omega_i=\sqrt{f_i/k_f}\quad(\text{as in M1.24--M1.25}) \tag{M2.18}\label{eq:m2.18}
\end{align}
\end{eqbox}

\begin{dimtable}
\vu_{\text{trim}} & hover input (operating point) & \R^4 & N, N\,m\\
\vu_k^{\text{abs}} & absolute thrust and torque & \R^4 & N, N\,m\\
\vOm & motor speeds & \R^4 & rad/s\\
\end{dimtable}

\subsubsection*{Derivation}

Models and controllers of Model~2 work in deviations \eqref{eq:m2.1}. The real motors need absolute values, so the
hover trim is added back. The rest is the same allocation as Model~1 (\cref{sec:m1b8}). The loop is closed by
measuring the new state of the \emph{nonlinear} drone (in the project: the EKF of Model~1) and starting block~9 again.
Logged data can also be fed back to block~1 to re-identify the model (outer dashed loop of \cref{fig:flow2}).

\begin{toybox}[one control sample, then a whole flight]
Drone at $[1.5,\,-1.0,\,0.8]$\,m, at rest; target $[0,0,2]$. The MPC gives $\delta\vu=[10.285,\,-0.279,\,-0.300,\,0.002]$, so
$\vu=[25.000,\,-0.279,\,-0.300,\,0.002]$, $\vf=[7.115,\,5.327,\,6.539,\,6.019]$\,N and
$\vOm=[912,\,789,\,875,\,839]$\,rad/s.

\textbf{Nonlinear flight}, 10\,s with gusts and sensor noise, controllers designed only on the DMDc model:
\begin{center}\small
\begin{tabular}{lrrr}
\toprule
controller & max $|v_i|$ [m/s] & error at 10\,s [m] & motor speeds [rad/s]\\
\midrule
MPC & 0.994 & 0.015 & 584 -- 913\\
LQR (clipped) & 1.118 & 0.031 & 613 -- 911\\
\bottomrule
\end{tabular}
\end{center}
Both stabilise the real drone. Only MPC honours \code{max\_velocity} = 1\,m/s.
\end{toybox}

\begin{figure}[h]
  \centering
  \includegraphics[width=0.95\linewidth]{step10_nonlinear_flight.png}
  \caption{Toy model, Model~2 step~10: the nonlinear drone flown by controllers that were designed on a model learned
  from data.}
\end{figure}

\begin{projbox}
This is the moment of truth: a controller designed on a learned, linear, approximate model flies the real, nonlinear,
gusty, motor-limited drone. In the project this block is the same motor mixer as Model~1, fed through ArduPilot.
\end{projbox}

\begin{codebox}
\intoy \code{step10\_apply\_to\_drone.m}, \code{lib/motor\_allocation.m}.
\end{codebox}

% ---------------------------------------------------------------------
\subsection{Recap of Model 2}
\label{sec:m2recap}
% ---------------------------------------------------------------------

\begin{center}\small
\begin{tabularx}{\linewidth}{@{}c >{\raggedright\arraybackslash}p{2.9cm} >{\raggedright\arraybackslash}p{3.2cm} >{\raggedright\arraybackslash}p{3.0cm} X@{}}
\toprule
\textbf{\#} & \textbf{Block} & \textbf{In} & \textbf{Out} & \textbf{The idea in one sentence}\\
\midrule
1 & Flight data & drone, excitation & $\{\vx_k\},\{\vu_k\}$ & Record what the drone did and what we told it, as deviations from hover.\\
2 & DMDc & $X$, $X^+$, $U$ & $A$, $B$ & Least squares finds the linear rule that best predicts the next sample.\\
3 & Bellman & $A,B,Q,R$ & $V(\vx)=\vx^\top P\vx$ & The best plan from here is the best first move plus the best plan from where it leads.\\
4 & LQR & $A,B,R,P$ & $K$ & Minimising a quadratic gives a linear law $\vu=-K\vx$.\\
5 & Riccati recursion & $P_0$ & $P_\infty$ & Look one more step ahead each iteration until nothing changes.\\
6 & DARE & $A,B,Q,R$ & $P$, final $K$ & Jump straight to that steady state; $\vx_0^\top P\vx_0$ is the optimal cost.\\
7 & MPC & model, limits, $P_f$ & optimisation problem & Plan the next second optimally without breaking any rule.\\
8 & QP & problem, $\vx_0$ & $H,\bm f,G,\bm h$, $\mathbf U^*$ & Substitute the predictions: a convex bowl inside a fence.\\
9 & MPC input & $\vx_k$ & $\vu_0^*$ & Apply only the first move, then re-plan with fresh data.\\
10 & Apply to drone & $\vu_0^*$ or $-K\vx_k$ & $\Omega_1..\Omega_4$ & Add hover thrust back, split among rotors, fly the real drone.\\
\bottomrule
\end{tabularx}
\end{center}
